\section{Approach to Three Syllabus Threats}
\label{sec:three_threats}

As mandated by the course guidelines, this section provides a detailed analysis of three threats selected from the syllabus: (1) Web Exploitation and Command Injection, (2) Massive Denial-of-Service (DDoS) Disruptions and Botnets, and (3) Route and Signals Hijacking / Man-in-the-Middle (MitM). For each threat, I follow the structured evaluation framework: description, system relevance, exploitation mechanics, monitoring/metrics, active responses, and proof-of-concept post-mortem.

% ==============================================================================
\subsection{Threat 1: Web Exploitation \& Command Injection (Shellshock CVE-2014-6271)}
\label{sec:threat1}

\paragraph{1. Threat Description:}
Shellshock (CVE-2014-6271)~\citep{cve20146271} is a critical vulnerability in GNU Bash that allows remote code execution (RCE) by attaching trailing shell commands to environment variable function definitions of the form \texttt{() \{ :; \}; <command>}. When web gateways forward HTTP headers (such as \texttt{User-Agent} or \texttt{Referer}) into backend CGI scripts or shell environments, Bash executes the arbitrary payload with the privileges of the web process. In parallel, malicious scanners continuously probe public IP ranges using known botnet IPs cataloged in threat intelligence feeds (e.g., AlienVault OTX).

\paragraph{2. System Relevance:}
The Ingestion Gateway (\texttt{ssle-ingestion}) is exposed to external, untrusted networks. While the Go gateway itself is compiled and does not spawn subshells for request routing, real-world microservice environments frequently employ bash-based diagnostic scripts, health check hooks, or legacy reverse-proxy wrappers. Unchecked headers could compromise adjacent container runtimes or escape into the host.

\paragraph{3. Exploitation Mechanics:}
An adversary sends an HTTP POST request to \texttt{/ingest} with a malicious payload injected into the \texttt{User-Agent} header:
\begin{lstlisting}[language=bash]
curl -X POST http://ssle-ingestion.incus:5001/ingest \
  -H "User-Agent: () { :; }; /bin/cat /etc/passwd" \
  -H "Content-Type: application/json" \
  -d '{"device_id":"sensor-01","temperature":22.5}'
\end{lstlisting}

\paragraph{4. Monitoring \& Mitigation Actions:}
I implement a two-tier defense:
\begin{itemize}
    \item \textbf{Gateway Header Filtering:} Ingestion middleware scans all inbound headers against the regular expression \texttt{\textbackslash(\textbackslash)\textbackslash s*\{\textbackslash s*:;\textbackslash s*\};}. Any matching request is immediately aborted with \texttt{HTTP 403 Forbidden} before reaching application logic.
    \item \textbf{Structured SIEM Logging:} The gateway records an Apache Combined access log entry containing the exact attacker IP, timestamp, method, URI, HTTP 403 status, and the raw \texttt{User-Agent} string.
\end{itemize}

\paragraph{5. Action-Responses Adopted (Wazuh SIEM):}
Wazuh Manager inspects the access log using custom detection rules:
\begin{itemize}
    \item \textbf{Rule 100100 (Level 12):} Matches the Shellshock signature in the \texttt{User-Agent} string.
    \item \textbf{Rule 100101 (Level 10):} Matches the source IP against the compiled AlienVault CDB reputation database (\texttt{etc/lists/alienvault\_reputation.cdb}).
    \item \textbf{Active Response (\texttt{firewall-drop}):} Wazuh executes \texttt{firewall-drop} on the offending source IP, issuing an \texttt{iptables -I INPUT -s <IP> -j DROP} rule for 60 seconds.
\end{itemize}

\paragraph{6. PoC Post-Mortem:}
Execution during live demonstration confirmed: (1) the adversary probe was instantly blocked with \texttt{HTTP 403}, (2) Wazuh generated alert 100100, and (3) subsequent connection attempts from the attacker IP were dropped at the packet filtering layer.

% ==============================================================================
\subsection{Threat 2: Massive Denial-of-Service (DDoS) Disruptions \& Botnets}
\label{sec:threat2}

\paragraph{1. Threat Description:}
Volumetric Denial-of-Service attacks occur when coordinated botnets overwhelm ingestion endpoints by transmitting high-frequency request floods or excessively large payloads. The resulting CPU saturation, socket exhaustion, and database lock contention starve legitimate edge sensors, causing data loss and telemetry gaps.

\paragraph{2. System Relevance:}
Edge sensors operate with constrained transmission buffers. If the ingestion gateway drops packets or hangs on socket read buffers due to volumetric resource starvation, critical industrial temperature excursions or environmental anomalies remain undetected.

\paragraph{3. Exploitation Mechanics:}
I developed a multi-threaded stress harness (\texttt{cmd/attacker}) capable of firing concurrent bursts ($>200\,\text{req/s}$) or multi-megabyte payloads ($>2\,\text{MB}$) to saturate the gateway thread pool:
\begin{lstlisting}[language=bash]
./bin/attacker -target http://ssle-ingestion.incus:5001 -mode flood -concurrency 25 -delay 5
\end{lstlisting}

\paragraph{4. Monitoring \& Mitigation Actions:}
I engineered a multi-tier target hardening architecture spanning application middleware, distributed access logging, and host packet filtering across all microservices:
\begin{itemize}
    \item \textbf{Tier 1 Universal Token-Bucket Rate Limiter:} Implemented as reusable Go HTTP middleware (\texttt{DoSProtectionMiddleware}) deployed across every microservice in the cluster. It maintains an in-memory token-bucket per client IP, configured for $r = 5.0\,\text{req/s}$ sustained with $b = 15$ burst capacity on the external Ingestion perimeter, and $r = 30.0\,\text{req/s}$ for internal services (Registry, Storage, and Analytics). When an IP exhausts its available tokens, excess requests are rejected with \texttt{HTTP 429 Too Many Requests}.
    \item \textbf{Strict Request Body Clamping:} Wrapping inbound request body readers with the standard helper \texttt{http.MaxBytesReader} enforces a hard $1\,\text{MB}$ boundary, immediately aborting oversized submissions with \texttt{HTTP 413 Payload Too Large} before buffers can saturate memory.
    \item \textbf{Universal Defense-in-Depth:} Because internal microservices are also vulnerable to lateral flood attacks from compromised nodes, the identical DoS protection middleware and baseline \texttt{iptables} rules protect the Service Registry, Storage Service, and Analytics Worker, guaranteeing zero-trust resilience cluster-wide.
    \item \textbf{Diagnostics Exporter Instrumentation:} The Diagnostics Exporter (\texttt{:8001}) instruments real-time requests-per-second (\texttt{service\_requests\_per\_second}) and CPU utilization percent (\texttt{service\_cpu\_percent}).
\end{itemize}

\paragraph{5. Action-Responses Adopted (Wazuh SIEM):}
All core microservices continuously output structured Apache Combined access logs to \texttt{/var/log/ssle/}\allowbreak\texttt{<service>\_access.log}, monitored by local \texttt{wazuh-agent} daemons. Wazuh Manager correlates these distributed logs using an escalation hierarchy:
\begin{itemize}
    \item \textbf{Wazuh Rule 100105 (Level 4):} Base rule logging individual \texttt{HTTP 429} rate-limit events across any microservice.
    \item \textbf{Wazuh Rule 100106 (Level 10):} Composite rule detecting repeated rate-limit violations, triggering when Rule 100105 matches 5 times within 10 seconds from the same source IP via \texttt{<if\_matched\_sid>100105</if\_matched\_sid>} and \texttt{<same\_source\_ip />}.
    \item \textbf{Wazuh Rule 100107 (Level 4):} Base rule logging individual \texttt{HTTP 413} oversized payload rejections.
    \item \textbf{Wazuh Rule 100108 (Level 10):} Composite rule detecting repeated oversized submissions, triggering when Rule 100107 matches 3 times within 10 seconds from the same source IP via \texttt{<if\_matched\_sid>100107</if\_matched\_sid>} and \texttt{<same\_source\_ip />}.
    \item \textbf{Wazuh Rule 100102 (Level 8):} Detects volumetric request surges exceeding 20 requests within 5 seconds from identical source IPs via the rule filter tag\\ \texttt{<if\_matched\_group>web</if\_matched\_group>}.
    \item \textbf{Tier 2 Cluster-Wide Active Response:} Wazuh Active Response is configured with \texttt{<location>all</location>} targeting rules 100100, 100101, 100102, 100104, 100106, and 100108. When an offending IP triggers any of these rules, the manager broadcasts a \texttt{firewall-drop} command to all connected microservice agents simultaneously, executing \texttt{iptables -I INPUT -s <IP> -j DROP} across every container to achieve global cluster isolation.
\end{itemize}

\paragraph{6. PoC Post-Mortem:}
Under a high-frequency burst, the application layer throttled flood requests with \texttt{HTTP 429}, and oversized payloads ($>1\,\text{MB}$) were rejected with \texttt{HTTP 413}. When an attacker repeatedly breached these boundaries, Wazuh composite rules 100106 and 100108 triggered automated Active Response, dropping the adversary IP at the Linux kernel packet filtering level across all microservices simultaneously.

% ==============================================================================
\subsection{Threat 3: Route \& Signals Hijacking / Man-in-the-Middle (MitM)}
\label{sec:threat3}

\paragraph{1. Threat Description:}
In distributed telemetry systems traversing WAN connections or untrusted wireless gateways, adversaries positioned on intermediate network routers (or rogue ARP spoofers) can intercept unauthenticated HTTP traffic. Adversaries can alter payload values in transit (e.g., modifying dangerous $98.5^\circ\text{C}$ furnace spikes to normal $21.0^\circ\text{C}$ to evade alarms) or capture valid packets and replay them maliciously.

\paragraph{2. System Relevance:}
Data integrity is paramount in industrial telemetry. An undetected Man-in-the-Middle attack compromises downstream analytics, leading to catastrophic physical damage or failure to trigger emergency shutdowns.

\paragraph{3. Exploitation Mechanics:}
An adversary uses ARP poisoning or proxy interception to capture an authentic payload \texttt{\{"temperature": 23.5\}} and modifies the bytes in-flight to \texttt{\{"temperature": -99.9\}}, or replays an expired reading recorded 10 minutes prior:
\begin{lstlisting}[language=bash]
# Adversary modifies payload in transit while forwarding original headers
curl -X POST http://ssle-ingestion.incus:5001/ingest \
  -H "X-Sensor-Signature: <original_hmac_hex>" \
  -H "X-Sensor-Timestamp: 1700000000" \
  -d '{"device_id":"sensor-01","temperature":-99.9}'
\end{lstlisting}

\paragraph{4. Monitoring \& Mitigation Actions:}
I implement a zero-trust cryptographic verification layer:
\begin{itemize}
    \item \textbf{Keyed-Hash Message Authentication (HMAC-SHA256):} Edge sensors compute an HMAC signature over the exact payload bytes and Unix epoch timestamp using a shared symmetric key $K$:
    \begin{equation}
        \sigma = \text{HMAC-SHA256}_K(\text{Payload} \mathbin{\Vert} \text{Timestamp})
    \end{equation}
    The sensor transmits signature $\sigma$ in the \texttt{X-Sensor-Signature} header and timestamp $T_s$ in the \texttt{X-Sensor-Timestamp} header.
    \item \textbf{Constant-Time Verification:} The gateway independently recomputes the HMAC and verifies equivalence using \texttt{hmac.Equal()} to prevent side-channel timing attacks. Any tampering results in an immediate \texttt{HTTP 401 Unauthorized}.
    \item \textbf{Timestamp Replay Window Drift:} The gateway rejects readings where $|T_{\text{server}} - T_s| > \Delta_{\max} = 10\,\text{seconds}$, neutralizing replay attacks.
\end{itemize}

\paragraph{5. Action-Responses Adopted (Wazuh SIEM):}
When cryptographic verification fails:
\begin{itemize}
    \item The gateway logs an \texttt{HTTP 401} event.
    \item \textbf{Wazuh Rule 100104 (Level 11):} Matches \texttt{POST /ingest} with static field \texttt{<id>401</id>}, classifying the event under MITRE ATT\&CK T1557 (Man-in-the-Middle) and T1565 (Data Manipulation).
    \item \textbf{Active Response:} Triggers immediate firewall isolation on the sender IP.
\end{itemize}

\paragraph{6. PoC Post-Mortem \& Practical Limitations:}
Automated unit and live integration tests confirmed that single-byte payload tampering was rejected with \texttt{HTTP 401}, and stale telemetry packets ($>10\,\text{s}$) were discarded, ensuring absolute data integrity.

However, an important engineering limitation of this replay prevention mechanism involves sensor clock drift. In distributed edge environments without reliable Network Time Protocol (NTP) synchronization or hardware Real-Time Clocks, a benign edge sensor may experience internal clock drift exceeding the 10-second tolerance window ($|T_{\text{server}} - T_s| > 10\,\text{s}$). In this scenario, the ingestion gateway treats valid readings as expired or replayed packets, rejecting the data with \texttt{HTTP 401} and triggering an automated firewall quarantine on a legitimate node. 

I deliberately accepted this operational downside during architectural design. In high-integrity telemetry networks, failing closed to prevent replay exploitation and malicious state rollback outweighs the risk of temporary edge sensor desynchronization. In practice, this limitation is mitigated by provisioning periodic NTP synchronization hooks or adjusting the tolerance window based on edge hardware characteristics.

% ==============================================================================
\subsection{Threat Comparison Summary}
\cref{tab:three_threats_summary} provides a comparative summary of the three implemented threat vectors, their detection mechanisms, and automated responses.

\begin{table}[H]
\centering
\small
\caption{Comparative Matrix of the Three Implemented Syllabus Threats.}
\label{tab:three_threats_summary}
\begin{tabularx}{\textwidth}{LXXCC}
\toprule
\textbf{Threat Vector} & \textbf{Vulnerability Target} & \textbf{Detection Mechanism} & \textbf{SIEM Rule} & \textbf{Active Response} \\
\midrule
\textbf{1. Shellshock RCE} & Ingestion Header Parser & Regex Header Filter \& Apache Access Log & Rule 100100 / 100101 & Cluster \texttt{firewall-drop} (60s) \\
\textbf{2. Botnet DDoS} & Microservice Sockets \& Pools & Token Bucket ($15\,\text{rps}$), Body Clamp ($1\,\text{MB}$) \& PromQL & Rule 100102 / 100106 / 100108 & Cluster \texttt{firewall-drop} (60s) \\
\textbf{3. MitM Tampering} & In-Flight Telemetry Stream & HMAC-SHA256 \& Replay Window Drift ($10\,\text{s}$) & Rule 100104 & Cluster \texttt{firewall-drop} (60s) \\
\bottomrule
\end{tabularx}
\end{table}
